\chapter{总结与展望}

\section{全文总结}

本文围绕细粒度猫狗品种识别任务,研究了基于 CLIP 的提示学习方法在少样本条件下的应用问题。针对 CoOp 使用静态提示、CoCoOp 未显式区分样本难度的不足,本文设计并实现了 DynamicPromptTrainer,将图像条件提示生成与难度感知样本加权结合到同一训练框架中。

在方法上,本文完成了以下几项工作。首先,构建 AdaptivePromptLearner,在基础上下文向量之外加入类别自适应因子,使不同类别能够使用不同的提示缩放方式。其次,设计 SoftPromptAdapter,利用两层 MLP 根据图像特征生成提示偏移,让提示词不再固定于单一模板。再次,引入 DifficultyWeightCalculator,根据真实类别预测概率和误分类情况对训练样本赋予不同权重,使困难样本在损失函数中的贡献更大。最后,在原有分类系统之外梳理并整合了推理时安全扩展模块,说明其基于 TTC 的检测与防御机制以及与动态提示模型的集成方式。

在实验上,本文以 Oxford-IIIT Pets 数据集为测试对象,完成了 1-shot、2-shot、4-shot、8-shot 和 16-shot 五组比较实验。结果显示,本文方法在 2-shot、4-shot、8-shot 和 16-shot 设置下均超过 CoCoOp,其中 16-shot 达到 89.8\% 的最高准确率。这说明在细粒度宠物分类任务中,对提示词进行图像条件化调整的同时,再引入面向困难样本的训练权重分配,能够带来稳定收益。

\section{工作不足}

虽然本文方法取得了较好的结果,但当前工作仍有几个明显限制。

第一,实验目前集中在 Oxford-IIIT Pets 单一数据集上,任务类型比较统一。该结果能够说明方法在宠物品种识别上的有效性,但还不足以证明其在鸟类、车型或植物等其他细粒度场景中的泛化能力。

第二,项目中已经完成主实验对比与图表整理,但尚未形成系统的模块级消融实验。例如,单独移除难度权重模块、固定类别自适应因子或替换不同提示初始化方式后的效果,还需要进一步补充。

第三,安全模块虽然已经完成代码集成,并提供了独立评估脚本,但当前本地环境尚未具备完整的数据与依赖条件,因此还没有形成可直接写入论文正文的最终防御实验结果。这意味着安全部分目前主要完成了系统设计与评估流程层面的整理。

\section{后续展望}

后续工作可以从以下几个方面继续推进。

\begin{enumerate}
\item 扩展数据集范围。在保持当前实现框架不变的前提下,将方法迁移到 CUB、Stanford Cars 等细粒度数据集,验证难度感知机制是否具有更普遍的效果。
\item 补充更完整的消融实验。围绕 SoftPromptAdapter、DifficultyWeightCalculator 和类别自适应因子分别开展控制变量实验,进一步解释性能提升来自哪些具体模块。
\item 完善安全评估环境。修正 \texttt{evaluate\_security.py} 中的数据集加载问题,补齐测试数据路径,在统一攻击强度下系统评估防御前后干净样本与对抗样本的准确率变化。
\item 尝试更强的主干网络与提示结构。例如使用 ViT 系列 CLIP 主干,或者探索更长上下文、分层提示和多模板提示的组合方式。
\end{enumerate}

\section{本章小结}

本文完成了一个面向细粒度猫狗分类的动态提示学习系统,并给出了可复现的实验结果。同时,项目还引入了推理阶段的安全扩展模块,为后续开展对抗鲁棒性研究预留了接口。当前工作已经验证了方法的基本有效性,但仍有扩展数据集、完善消融实验和补全安全评估等进一步研究空间。
